\chapter{Introdução}
\label{cap:01}

O objetivo deste documento é esclarecer aos autores o formato que deve ser utilizado nos relatórios técnicos a serem submetidos ao final dos cursos de Graduação e Pós-Graduação do IFSP câmpus São João da Boa Vista. Este documento está escrito de acordo com o modelo indicado para a formatação dos relatórios técnicos; assim, serve de referência, ao mesmo tempo em que comenta os diversos aspectos da formatação.

Observe as instruções e formate seu relatório técnico de acordo com este padrão. Lembre-se que uma formatação correta contribui para uma boa avaliação do seu trabalho.

Além disso, neste documento estão listadas as seções obrigatórias que você deverá fornecer, bem como os exemplos dos comandos mais comuns que serão utilizados na construção de seu documento. Para pesquisar sobre mais comandos, recomenda-se a utilização do site \url{https://ctan.org/}, que é a biblioteca principal do \LaTeX, e o do site \url{https://tex.stackexchange.com} que é uma das principais comunidades para solução de dúvidas relacionadas a \LaTeX. Ambas são em inglês.

A introdução é um elemento preliminar, opcional, utilizado para fornecer informações específicas, comentar tecnicamente o conteúdo do trabalho, além de evidenciar as motivações que levaram o autor à escolha de determinado tema.

Trata-se de importante estratégia de aproximação, pois permite valorizar a escolha do assunto, mostrar a relevância da abordagem temática e esclarecer quanto ao passo-a-passo utilizado na estruturação do texto.

Na introdução, o leitor terá condições de avaliar:

\begin{itemize}
	\item O grau de informação, conhecimento e competência técnica do autor relativamente ao assunto a ser tratado;
	\item A qualidade, a eficiência, a originalidade e o ineditismo de sua abordagem;
	\item A pertinência das informações apresentadas e a possibilidade de acrescentar algo de novo ao universo conceitual do leitor.
\end{itemize}


% Substitui, em capitulo01Introducao.tex, as seções "Objetivos", "Objetivo Geral" e "Objetivos Específicos" do template.
% Usa enumitem (já carregado em estrutura.sty). Os rótulos oe:* são referenciados na Metodologia.

\section{Objetivos}
\label{sec:objetivos}

\subsection{Objetivo Geral}

Desenvolver um protótipo de ferramenta computacional para a análise comparativa de proteínas associadas a vias metabólicas do sabor e do aroma do café entre variedades de \textit{Coffea arabica} e espécies relacionadas do gênero \textit{Coffea}.

\subsection{Objetivos Específicos}

\begin{enumerate}[label=OE\arabic*., ref=OE\arabic*]
	\item \label{oe:dados} Obter os proteomas preditos das variedades e espécies de \textit{Coffea} que serão comparadas e verificar sua qualidade por critérios objetivos;
	\item \label{oe:vias} Definir as vias metabólicas associadas ao sabor e ao aroma e a sequência-âncora de referência de cada via;
	\item \label{oe:localizacao} Implementar a localização do gene correspondente à sequência-âncora em cada proteoma;
	\item \label{oe:alinhamento} Implementar o alinhamento global com penalidade de lacuna afim e a métrica de similaridade entre pares de proteínas;
	\item \label{oe:modulo} Organizar o núcleo bioinformático em um módulo reutilizável que gere os resultados em formato estruturado;
	\item \label{oe:software} Desenvolver uma API e uma interface web para executar e visualizar as comparações;
	\item \label{oe:validacao} Validar a correspondência dos genes encontrados e a métrica de similaridade com dados da literatura.
\end{enumerate}
